
\section{主要风险}
\begin{riskbox}[七项核心风险]
\begin{enumerate}
\item 光模块利润集中度高，单季高增长不能自动外推为全年和长期。
\item USD1.2bn扩产、AI PCB扩产与并购贷款增加折旧、营运资金和财务费用。
\item 客户认证、订单金额、ASP、利用率和良率缺口可能导致收入和利润预测下修。
\item EML、硅光、CPO和铜互联技术路线变化带来替代风险。
\item 商誉和并购整合风险：Q1商誉约47.69亿元，其中索尔思并购商誉约28.25亿元。
\item 传统FPC、显示和汽车精密组件仍受消费电子、汽车和出口周期影响。
\item 2025年外销收入占比81.41\%，汇率、贸易政策和海外经营可能影响利润。
\end{enumerate}
\end{riskbox}

\section{催化剂与失效}
\begin{exhibitbox}[表7-1\quad 监控清单]
\begin{tabularx}{\exhibitboxwidth}{L{3.0cm}X L{3.2cm}X}
\textbf{方向} & \textbf{正向催化剂} & \textbf{失效信号} & \textbf{估值动作} \\
H1/正式半年报 & 预告归母29-30亿元、扣非24-25亿元；正式半年报需确认 & 正式数据明显低于预告或光模块分部仍无改善 & 保留公司层信用，降低分部信用 \\
光模块 & 收入超过24亿元、分部毛利和现金流可核验 & 仅有扩产新闻，收入/毛利不跟随 & 扩产只保留期权价值 \\
AI PCB & 高端产品收入、ASP、利用率和良率披露 & 电子电路仍只有总额披露 & 不提高AI PCB倍数 \\
财务质量 & CFO/NP改善，融资成本稳定 & 借款和财务费用上升而利润不增 & 提高保守准备 \\
价格行为 & 回撤至支持区后有基本面确认 & 跌破221.68且业绩转弱 & 降低风险预算 \\
\end{tabularx}
\sourcenote{analysis/risk\_framework.md；analysis/implied\_growth\_sensitivity.md}
\end{exhibitbox}

\section{最终行动建议}
对于已有仓位，建议把持仓逻辑从“题材重估”切换为“业绩验证”：保留能够承受高波动的核心观察仓，避免在涨停和天量成交时追高。对于新增仓位，优先等待价格回到支持区，或等待正式半年报验证光模块收入、毛利率和现金转换。若验证通过，基准目标可上修；若验证失败，基本面锚和市场锚同时下移。
